\documentclass[11pt]{article}

\usepackage[utf8]{inputenc}
\usepackage[T1]{fontenc}
\usepackage[margin=1in]{geometry}
\usepackage{enumitem}
\usepackage{hyperref}
\usepackage{parskip}
% \usepackage{microtype}

\setlist[enumerate,1]{label=(\arabic*), leftmargin=*, itemsep=0.5em}
\setlist[enumerate,2]{label=(\alph*), leftmargin=2em, itemsep=0.35em}
\setlist[enumerate,3]{label=(\roman*), leftmargin=2em, itemsep=0.25em}
\setlist[itemize]{leftmargin=2em, itemsep=0.35em}

\newcommand{\articleheading}[2]{%
    \section*{Article #1. #2}
    \addcontentsline{toc}{section}{Article #1. #2}
}

\begin{document}

\begin{center}
    {\LARGE\bfseries WPI IEEE Student Branch Constitution}
\end{center}

\tableofcontents
\newpage

\articleheading{I}{Name}

\begin{enumerate}
    \item The name of this organization will be the Worcester Polytechnic Institute Institute of Electrical and Electronics Engineers Student Branch and hereafter referred to as WPI IEEE or as the Branch.
\end{enumerate}

\articleheading{II}{Purpose}

\begin{enumerate}
    \item The purpose of the Branch shall be to:
    \begin{enumerate}[label=\Roman*.]
        \item Represent the National Institute of Electrical and Electronics Engineers (IEEE) organization at Worcester Polytechnic Institute (WPI).
        \item Provide technical, professional, and social events for the WPI community.
    \end{enumerate}
\end{enumerate}

\articleheading{III}{Membership}

\begin{enumerate}
    \item There will be three classes of membership in the Branch:
    \begin{enumerate}[label=\Roman*.]
        \item \textbf{General Member}
        \begin{enumerate}
            \item A qualified member of the community becomes a general member if they attend one event and register on MyWPI.
        \end{enumerate}

        \item \textbf{Active Member}
        \begin{enumerate}
            \item Active member status shall be granted to WPI students who have at least 3 membership points. Membership points are earned in the following ways:
            \begin{enumerate}
                \item 1 point is earned for attending an event.
                \item 0.5 points are earned for attending an executive board meeting.
            \end{enumerate}
            \item Membership points are accumulated across a rolling semester. A rolling semester is defined as a contiguous span of two terms where a student is not on leave. If a student is on leave, they must notify the executive board to exclude the term they are away from the rolling period.
            \item Under extenuating circumstances, the Executive Board may waive the above requirements.
            \item In order to obtain and maintain active membership, the member must be in good academic standing as defined by WPI.
        \end{enumerate}

        \item \textbf{Away}
        \begin{enumerate}
            \item Active members who leave WPI for a period of off-campus academic activity will automatically obtain away member status. Members are responsible for informing the Executive Board that they will be away from campus. Away members are not expected to fulfill the meeting requirements necessary to maintain active status.
        \end{enumerate}

        \item \textbf{Executive}
        \begin{enumerate}
            \item Any active member, except the Advisor, who serves as a member of the Branch's Executive Board, as defined in Article IV.
        \end{enumerate}

        \item \textbf{Interim Executive}
        \begin{enumerate}
            \item An executive who is elected only for a specific time period to fulfill an executive position while the position is temporarily vacant.
        \end{enumerate}
    \end{enumerate}

    \setcounter{enumi}{2}
    \item The Branch will not discriminate based on race, gender, creed, religion, color, age, sexual orientation, disability, veteran status, or national origin.
    \item The Branch and its membership will not practice any physically or psychologically abusive hazing either intentionally or unintentionally.
\end{enumerate}

\articleheading{IV}{Executive Board}

The executive board shall serve as the central authority of the Branch. The executive board shall be composed of the following positions, hereafter referred to as executives:

\begin{enumerate}[label=\arabic*.]
    \item Chair
    \item Vice Chair
    \item Treasurer
    \item Secretary
    \item Events Chair
    \item Projects Chair
    \item Public Relations Chair
    \item Web Chair
\end{enumerate}

Executives shall be expected to:

\begin{itemize}
    \item Maintain the status of an active member.
    \item Be in attendance, excused, or on leave for all Branch activities, including executive board meetings.
    \item Satisfy all requirements specific to the executive's position.
    \item Share the responsibilities of the executive board.
    \item Maintain familiarity with the Branch bylaws.
    \item Maintain IEEE Student Membership throughout the duration of their term as an executive.
\end{itemize}

\subsection*{(1) Chair}

The Chair may also be referred to as the President in accordance with WPI Executive Board regulations. It shall be the responsibility of the Chair to oversee the operations of the Branch in its entirety. This duty shall entail:

\begin{itemize}
    \item Meeting with executives to discuss and develop goals and achievements.
    \item Moderating executive board meetings.
    \item Assigning tasks to executives that fall outside their designated responsibilities.
    \item Preparing an annual Branch plan during A-Term and submitting it in vTools.
    \item Submitting at least four event reports in vTools over the academic year.
    \item Maintaining executive reporting in vTools.
    \item Serving as a representative on the Department of Electrical and Computer Engineering (ECE) Student Advisory Board under the ECE Department Head.
    \item Maintaining the Branch Bylaws.
\end{itemize}

\subsection*{(2) Vice Chair}

The Vice Chair may also be referred to as the Vice President in accordance with WPI Executive Board regulations. It shall be the responsibility of the Vice Chair to oversee the internal operations of the Branch. This duty shall entail:

\begin{itemize}
    \item Assisting the Chair with their duties.
    \item Helping fellow executives discuss and develop goals and achievements.
    \item Holding fellow executives accountable for failing to meet responsibilities.
    \item Providing fellow executives with resources and assistance as needed.
    \item Upholding the Bylaws.
    \item Ensuring strong coordination and cooperation amongst all executives.
\end{itemize}

\subsection*{(3) Treasurer}

It shall be the responsibility of the treasurer to oversee the Branch finances. This duty shall entail:

\begin{itemize}
    \item Serving as the financial advisor for the Branch.
    \item Submitting an annual budget proposal to the WPI Student Government Association.
    \item Submitting Funding Requests to the WPI Student Government Association a minimum of 2 weeks before the intended use date.
    \item Pursuing funds from external sources including, but not limited to, the IEEE Worcester County Section (WCS), the WPI ECE department, and corporate sponsors.
    \item Maintaining Branch funds held by the ECE department.
    \item Maintaining a history of Branch finances.
    \item Completing reimbursements for pre-approved purchases.
\end{itemize}

\subsection*{(4) Secretary}

It shall be the responsibility of the secretary to oversee the administration of the Branch. This duty shall entail:

\begin{itemize}
    \item Recording meeting minutes during executive board meetings and sending out action items afterward.
    \item Maintaining member attendance records at Branch activities.
    \item Updating member activity status as required.
    \item Maintaining the Branch file system, which includes the Google Drive folder.
    \item Creating and sending out the Branch weekly newsletter.
    \item Ensuring synchronization between myWPI members and the email alias.
\end{itemize}

\subsection*{(5) Events Chairs}

There shall be two Events Chairs, and it will be the responsibility of the Events Chairs to coordinate all logistics and planning for Branch activities, including coordinating with other executives when needed. This duty shall entail:

\begin{itemize}
    \item Arranging weekly events, which must include the following each term of the academic year:
    \begin{itemize}
        \item One General Body Meeting (GBM)
        \item One technical event
        \item One social event
        \item One flagship event
    \end{itemize}
    \item Uploading event details into MyWPI 1.5--2 weeks before the event takes place.
    \item Filling out the event template 1.5--2 weeks before said event takes place.
    \item Being the main point of contact for other student groups on campus for collaboration.
    \item Maintaining the events calendar.
\end{itemize}

\subsection*{(6) Projects Chair}

It will be the responsibility of the Projects Chair to create and maintain technical projects that promote and/or benefit the Branch and WPI community. This duty shall entail:

\begin{itemize}
    \item Identifying at least one project that can be completed or improved upon every academic year.
    \item Maintaining existing projects to ensure that they are operational.
    \item Open-sourcing and maintaining all projects on the branch's GitHub organization.
    \item Coordinating and managing the Projects Committee.
    \item Managing the Projects budget.
\end{itemize}

\subsection*{(7) Public Relations Chair}

It will be the responsibility of the Public Relations Chair to promote and document Branch events. This duty shall entail:

\begin{itemize}
    \item Designing graphics for events.
    \item Promoting events and meetings in ways including, but not limited to, posting to social media, hanging posters and fliers, and advertising on campus televisions.
    \item Documenting events by taking photos and/or videos and posting recaps.
    \item Coming up with ideas for WPI IEEE Branded Merchandise.
\end{itemize}

\subsection*{(8) Web Chair}

It will be the responsibility of the Web Chair to maintain, update, and improve the Branch website. This duty shall entail:

\begin{itemize}
    \item Ensuring the website is publicly accessible at all times.
    \item Ensuring the website's code base is organized and version controlled.
    \item Making sure events, executives, projects, and other content on the website is up to date.
    \item Maintaining the website calendar.
    \item Improving the website by adding or improving features each academic term.
\end{itemize}

\articleheading{V}{Meetings}

The Branch shall have four classes of meetings:

\begin{enumerate}
    \item \textbf{General Body Meeting:} A meeting held at least once per term, ideally at the beginning of the term, that is open to all WPI students. During general body meetings, chapter updates and future events shall be discussed. Elections may also be run during or in place of the C-Term General Body Meeting.

    \item \textbf{Events:} Organized activities, gatherings, or initiatives held to foster engagement, promote the club's mission, and provide opportunities for member involvement. Events may include workshops, social gatherings, professional development opportunities, and other activities aligned with the club's objectives.

    \item \textbf{Flagship Event:} Flagship events are large-scale events that happen once per academic term. The current set of flagship events is described below. Each flagship event requires the formation of a committee (see Branch Bylaws for descriptions of committees) to assist in the planning and execution of these large events.
    \begin{enumerate}
        \item \textbf{A-Term: PCB Design Class:} A project-oriented hardware design course where students learn PCB design over several weeks, then design, receive, and assemble their own PCBs. (Location: Flexible, Date: Early September)
        \item \textbf{B-Term: Spark Party:} A night to celebrate the ECE department's accomplishments, featuring music and dance performances from student groups, WCS demonstrations, faculty and MQP presentations, with a finale Tesla coil performance. (Location: AK 116, Date: Early December)
        \item \textbf{C-Term: Networking Night:} A unique opportunity for highly motivated, talented, technical WPI students across different majors to connect with top companies in the industry over a catered dinner. (Location: CC Odeum, Date: Late January)
        \item \textbf{D-Term: Hackathon:} A weekend-long project-focused hackathon with ample hardware and facilities open to students over the weekend. (Location: Unity Hall, Date: Late March)
    \end{enumerate}

    \item \textbf{Executive Board Meeting:} A 1-hour meeting held once per week during all academic terms that is open to the WPI community. During executive board meetings, the executive board discusses and plans Branch activities and votes on Branch matters.
\end{enumerate}

\articleheading{VI}{Elections \& Officer Replacement/Removal}

\begin{enumerate}
    \item \textbf{Schedule:} The election of the executive board, hereafter referred to as the election, shall be held at least once a year and shall take place no later than the second-to-last full week of C term.

    \item \textbf{Nominations:} The executive board shall notify all active members of the election two weeks prior to the election. Any active member may nominate any undergraduate active member. Active members may be nominated for more than one position. Nominations may be submitted to the executive board up until the start of the election. Upon receiving a nomination, the executive board will immediately notify the nominee. To be a candidate in the election, the nominee must accept the nomination prior to the start of the election. The nomination is assumed to be automatically accepted in the case of a self-nomination.

    \item \textbf{Voting Membership:} The voting membership is the set of people who are eligible to vote in a certain ballot. For executive elections, the voting membership is the set of active members.

    \item \textbf{Moderators:} Elections for each board position shall be moderated by the two highest-ranking board members who do not intend to run for any elected position. In the event that no such members are available, the current Executive Board shall oversee the election of the Chair. Once a Chair has been elected by the general body, the newly elected Chair shall assume responsibility for moderating the elections for the remaining board positions.

    \item \textbf{Procedure:} When electing a board member, the following procedure will be followed:
    \begin{enumerate}
        \item The names of all candidates for the position shall be read.
        \item All candidates shall be asked to leave the room.
        \item Candidates shall be selected one by one, in a random order. Each selected candidate shall return to the room and may deliver a speech no longer than three minutes. The general body may then ask the selected candidate questions for up to five minutes, following which the selected candidate shall be asked to leave the room.
        \item Following this, the general body may discuss the candidates for up to 10 minutes.
        \item All candidates shall return to the room, and balloting to elect an executive shall occur as described below.
    \end{enumerate}

    \item \textbf{Balloting:} Balloting to elect an executive shall occur during a general body meeting and shall proceed as follows:
    \begin{enumerate}
        \item If there are no candidates for the position, the position shall be left vacant until an appointment can be made or the duties of the unfilled position have been equitably distributed among the other executives.
        \item If there are two candidates for the position, each member of the voting membership shall be entitled to an anonymous vote for one candidate. The candidate with the largest number of votes shall be elected to the executive position. In the case of a tie, current executive members shall hold a ballot to determine the winner. Executive members cannot participate in this ballot if they are one of the tied candidates. If no consensus is reached during this ballot, the current Chair will determine the new executive member.
        \item In the event that more than two candidates are nominated for a single position, the election shall be conducted using an anonymous, ranked-choice ballot. Every member of the voting membership shall be given the right to cast a single ballot, on which they must rank up to three candidates in order of their preference (i.e., first choice, second choice, and third choice).
    \end{enumerate}

    \item \textbf{Executive Vacancy:} While a board position remains vacant, the Executive Board shall determine an appropriate method for distributing the associated responsibilities among current board members until the position is formally filled.

    \item \textbf{Interim Appointments:} During any period in which a board position is vacant, the Executive Board shall make reasonable efforts to appoint an interim executive on a term-by-term basis. This process shall include distributing an application to all active members at the beginning of the term and selecting a suitable candidate for appointment. In the event that an interim is appointed to fill a vacant board position, the interim shall assume all responsibilities and duties associated with that position, as outlined in Article IV. These responsibilities shall be upheld until the original executive returns from leave or until the next general board elections are conducted.

    \item \textbf{Executive Removal:} An executive may only be removed from office when said executive becomes inactive, does not meet membership standards, or violates WPI's, local, state, or federal laws. Also, an executive may be removed by a two-thirds vote of the rest of the executive board. All cases and reasons for removal from office must be approved by the advisor before removal. Students can approach the advisor with concerns they may have about an executive.
\end{enumerate}

\articleheading{VII}{Finances}

\begin{enumerate}
    \item The Branch will primarily be funded through its own fundraising efforts, WPI Student Government Association and through national or regional IEEE chapters at their discretion.
\end{enumerate}

\articleheading{VIII}{Advisor}

\begin{enumerate}
    \item The Faculty Advisor must be an active IEEE member and is appointed by the IEEE Regional Director upon the recommendation of the Student Members of the Branch, the Regional Student Activities Committee Chairman, and the Electrical \& Computer Engineering Department of WPI. The appointment or reappointment shall normally be for two years and begin on July 1st. The advisor is a non-voting member of the Branch.
\end{enumerate}

\articleheading{IX}{Amendments/Revisions}

\begin{enumerate}
    \item \textbf{Schedule:} In accordance with WPI SAO regulations, the Branch's constitution shall be reviewed and updated once a year.

    \item \textbf{Review Process:} To update the Branch's constitution, the Chair and Vice Chair must meet to discuss and draft possible changes. When all the desired updates have been drafted, the Chair and Vice Chair will present these changes to the executive board at an executive board meeting. Each change, individually, will then be discussed and voted on.
\end{enumerate}

\articleheading{X}{SOC Enabling Clause}

\begin{enumerate}
    \item The WPI IEEE Student Branch agrees to abide by the policies of Worcester Polytechnic Institute as well as all federal, state, and local laws.
\end{enumerate}

\end{document}
